% =====================================================================
% Results section ("5. Numerical Experiments") for
% neural_physics_uav_flight_path_optimisation
%
% Integration notes:
%   * \input this file BEFORE the current "5. Discussion" section;
%     Discussion and Conclusion then renumber to 6 and 7 automatically.
%   * Requires: graphicx, booktabs, amsmath, amssymb (all standard
%     with elsarticle).
%   * Figures live next to this file; add to the preamble:
%       \graphicspath{{Q1/figs/}{Q2/figs/}{Q3/figs/}{Q4/figs/}}
%     (or copy the eight fig0X_*.pdf files into the main tex folder).
%   * Every figure/table is regenerable from the repository:
%     experiments/exp*.py -> output/runs/<exp>/ -> tools/visualize/
%     paper/fig0X_*.py (one script per figure, pinned to git SHAs).
%   * The two appendix subsections at the end of this file
%     (\paragraph blocks marked APPENDIX) should be moved into the
%     paper's appendix if one exists.
% =====================================================================

\section{Numerical experiments}
\label{sec:experiments}

This section validates the wave--advection routing framework of
Sections~2--4 in four steps, each addressing one question a
practitioner would ask. First, \emph{is the computed arrival time
correct?} We verify the solver against analytic solutions in free
space, uniform wind, anisotropic climb/descent, and single-obstacle
diffraction (\S\ref{sec:verification}). Second, \emph{how close is it
to the discrete optimum on a real city?} We compare against an
anisotropic Dijkstra ground truth on a $118$-station district at
$5\,$m resolution (\S\ref{sec:groundtruth}). Third, \emph{what does
the PDE formulation buy?} We quantify the adjoint and superposition
shortcuts and the raw GPU throughput (\S\ref{sec:efficiency}).
Fourth, \emph{where does it break?} We measure the wavelength
operating envelope and the behaviour under spatially varying,
non-solenoidal wind (\S\ref{sec:robustness}), and close with a
scale-transfer demonstration from a delivery multirotor to a regional
aircraft (\S\ref{sec:transfer}). We additionally qualify positional
uncertainty in a controlled zero-wind obstacle test
(\S\ref{sec:uncertainty-results}), then test concurrent UAV execution
and a continuously time-varying velocity field
(\S\ref{sec:dynamic-operations}).

% ---------------------------------------------------------------------
\subsection{Experimental setup}
\label{sec:setup}

\begin{table}[t]
\centering
\caption{Computational domains used throughout this section.}
\label{tab:domains}
\small
\begin{tabular}{lllll}
\toprule
Domain & Grid & Resolution & Cells & Purpose \\
\midrule
Analytic (empty / 1 box) & up to $480^3$ & $1.25$--$20\,$m & $\le 1.1\times10^8$ & exact-solution benchmarks \\
Nanshan & $1333\times1443\times42$ & $5\,$m & $8.1\times10^7$ & real city, 118 stations \\
Nanshan coarse & --- & $50\,$m & $2.9\times10^5$ & wind sweeps, ablations \\
Regional & $200\times200\times20$ & $400\times400\times600\,$m & $8\times10^5$ & scale transfer (aircraft) \\
\bottomrule
\end{tabular}
\end{table}

Table~\ref{tab:domains} lists the four domains. The Nanshan domain is
a $6.7\times7.2\,$km district of Shenzhen voxelised from building
footprints, with $118$ delivery stations; the $50\,$m coarsening of
the same district is used where parameter sweeps would be too costly
at $5\,$m. Unless stated otherwise all runs use the production
configuration: horizontal speed $c_0 = 15\,$m/s, climb/descent speeds
$c_z^{\uparrow}=6$, $c_z^{\downarrow}=4\,$m/s (encoded by the vertical
advection offset of Section~2.1.1), $\Delta t = 0.1\,$s, $4500$ time
steps, a Gaussian source of width $\sigma = 1.5$ cells and pulse
width $\tau = 0.4\,$s, an eight-layer absorbing sponge, and
single-precision arithmetic on one NVIDIA A100.

Four metrics are used consistently. The \emph{optimality gap} is
$(T_{\mathrm{method}} - T_{\mathrm{ref}})/T_{\mathrm{ref}}$ evaluated
on \emph{valid pairs}, defined as station pairs separated by
$700$--$3000\,$m (below $700\,$m the source near field dominates;
beyond $3000\,$m sponge proximity does). The \emph{path time} is the
metric-integrated travel time along a traced path --- never the raw
field value $T$ at the endpoint. \emph{Success} means a
$-\nabla T$ traced path reaches its destination within two cells
without entering a building. \emph{Runtime} is JIT-warm wall-clock,
median of three, on one A100.

First-peak detection introduces a constant arrival bias approximately
equal to the source pulse peak time $t_0 = 2.5\tau$. We calibrate it
once per pulse width by solving an empty companion domain and fitting
$T$ against distance: the slope gives the unbiased effective speed
and the intercept the bias ($0.945\,$s at the production point),
which is subtracted from all arrival times reported below.

% ---------------------------------------------------------------------
\subsection{Verification and accuracy}
\label{sec:verification}

\begin{figure}[t]
\centering
\includegraphics[width=\linewidth]{fig01_convergence}
\caption{Free-space accuracy. (a)~Relative arrival-time error versus
grid spacing for four pulse widths; (b)~measured arrival bias versus
the source peak time $t_0$, confirming the linear calibration model.
Errors are evaluated on journeys of $170$--$250\,$m, corresponding to
exact travel times of $11.3$--$16.7\,$s at $c_0=15\,$m/s.}
\label{fig:convergence}
\end{figure}

\begin{figure}[t]
\centering
\includegraphics[width=\linewidth]{fig02_wind_aniso}
\caption{(a)~Uniform-wind validation: arrival-time error by direction
sector for Mach numbers $q/c_0 \in \{0.1, 0.3, 0.5\}$.
(b)~Climb/descent anisotropy: arrival time along the vertical axis
versus the exact $d/c_z^{\uparrow}$ and $d/c_z^{\downarrow}$ lines.}
\label{fig:windaniso}
\end{figure}

\begin{figure}[t]
\centering
\includegraphics[width=0.85\linewidth]{fig03_obstacle}
\caption{Single-obstacle benchmark: first-arrival field and traced
path versus the exact taut-string geodesic over the box. The exact
journey time in panel~(a) is $12.87\,$s. Across the four cell-centred
grids, the resolution-dependent reference times span
$12.80$--$13.27\,$s and the calibrated numerical times span
$12.75$--$13.32\,$s. Panel~(b) reports the signed relative gap after
removal of the independently measured source-pulse bias.}
\label{fig:obstacle}
\end{figure}

\paragraph{Free space (Fig.~\ref{fig:convergence}).}
Against the exact $T^\ast = \lVert x - x_0\rVert / c_0$, the RMS
pointwise error at the production operating point ($\Delta x = 5\,$m,
$\tau = 0.4\,$s) is $0.41\%$, decreasing monotonically to
$0.13$--$0.27\%$ at $\Delta x = 2.5\,$m; the fitted propagation speed
is accurate to better than $1.2\%$ at $5\,$m and $0.6\%$ at $2.5\,$m.
We deliberately do \emph{not} claim the nominal second-order
convergence of the spatial discretisation: the observable is the
first peak of a broadband pulse (subject to peak detection and
direction-dependent numerical dispersion), not an $L_2$ norm of a
smooth field, and the fitted order of the pointwise error is
$0.6$--$0.9$. The quantities that matter for routing --- propagation
speed and the calibrated bias --- converge much faster, as the
geodesic benchmark below confirms.

\paragraph{Uniform wind (Fig.~\ref{fig:windaniso}a).}
With a uniform background flow the analytic arrival time on every
ray is $d/(c_0 \pm q\cdot\hat d)$. For Mach numbers
$q/c_0 \le 0.3$ the median sector error is $\le 1.2\%$ (p95
$\le 3.3\%$) in all directions. At $M = 0.5$ the upwind sector
degrades to a median of $27\%$: the upwind wavelength
$(c_0 - q)\tau = 3\,$m falls below the grid spacing and the wave is
under-resolved. This is a \emph{measured} validity boundary of the
discretisation --- the remedy is a longer pulse or a finer grid ---
while the downwind and crosswind sectors at the same Mach number
remain at $3.1\%$ and $0.6\%$.

\paragraph{Climb/descent anisotropy (Fig.~\ref{fig:windaniso}b).}
The asymmetric vertical speeds are encoded exactly: fitted axis
speeds are $5.91\,$m/s up and $4.00\,$m/s down at
$\Delta x = 2.5\,$m (targets $6$ and $4$; errors $1.5\%$ and
$0.1\%$).

\paragraph{Single obstacle (Fig.~\ref{fig:obstacle}).}
Against the exact taut-string geodesic over a rectangular box, the
calibrated arrival-time gap is $\le 0.52\%$ across all resolutions
$\Delta x \in \{10, 5, 2.5, 1.25\}\,$m, and the $-\nabla T$ traced
path is homotopic to and geometrically close to the geodesic.
Diffraction and shadow zones are handled correctly --- the regime
that defeats straight-line heuristics.

\subsubsection{Uncertainty qualification without prescribed wind}
\label{sec:uncertainty-results}

\begin{figure}[t]
\centering
\includegraphics[width=\linewidth]{fig10_uncertainty}
\caption{Controlled positional-uncertainty experiment on the
single-obstacle geometry at $5\,$m resolution and zero prescribed wind.
(a)~Nominal paths obtained for constant
$D_{\mathrm{diff}}\in\{0,0.5,1,2.5,5,10,20\}\,\mathrm{m^2\,s^{-1}}$;
(b)~nominal obstacle clearance minus the exact three-dimensional
$95\%$ radial confidence radius along each path, where negative values
indicate that the uncertainty tube intersects the obstacle;
(c)~Monte Carlo collision probability from $4000$ trajectories per
diffusivity, with Wilson $95\%$ confidence intervals. The nominal
journey time remains $13.06$--$13.11\,$s across the sweep.}
\label{fig:uncertainty}
\end{figure}

We isolate uncertainty from wind by setting the prescribed velocity
field to zero and repeating the single-box benchmark of
Fig.~\ref{fig:obstacle}. For each constant $D_{\mathrm{diff}}$, an empty
companion solve removes the source-peak timing bias, after which the
arrival-time field is traced and $4000$ independent Brownian
realisations are sampled around the nominal trajectory. The exact
three-dimensional confidence radius is
$r_{95}=\sqrt{2\tau\chi^2_{3,0.95}}$, with
$\tau(t)=\int_0^t D_{\mathrm{diff}}\,\mathrm{d}t$. The sampled terminal
error p95 agrees with this analytical radius to within $2.8\%$ at every
non-zero diffusivity, providing an internal calibration check on the
uncertainty sampler.

The nominal route is insensitive to constant diffusivity at this
resolution: its travel time varies by less than $0.4\%$ and its minimum
clearance remains between $3.38$ and $4.31\,$m, with no monotonic
clearance increase (Fig.~\ref{fig:uncertainty}a). The stochastic safety
metrics, however, change strongly. At
$D_{\mathrm{diff}}=0.5\,\mathrm{m^2\,s^{-1}}$, $r_{95}=10.11\,$m and
the estimated collision probability is $17.75\%$ (Wilson $95\%$ CI
$16.60$--$18.96\%$); at $20\,\mathrm{m^2\,s^{-1}}$ these increase to
$63.98\,$m and $66.10\%$ ($64.62$--$67.55\%$), respectively. Thus the
diffusion model provides a quantitative risk envelope, but smoothing
of the wave field alone does not produce a monotonic geometric safety
margin in this benchmark. A clearance or chance constraint must be
applied to the uncertainty envelope if a prescribed collision-risk
tolerance is required.

\begin{figure}[t]
\centering
\includegraphics[width=\linewidth]{fig11_diffusion_resolution}
\caption{Grid sensitivity of constant uncertainty diffusion in the
zero-wind single-obstacle benchmark. (a)~Shift of the fitted empty-domain
peak-propagation speed relative to $D_{\mathrm{diff}}=0$ on each grid;
(b)~minimum nominal path clearance; (c)~nominal path time. The red cross
marks a tracer failure at $\Delta x=2.5\,$m and
$D_{\mathrm{diff}}=20\,\mathrm{m^2\,s^{-1}}$, where the explicit
diffusion number is $N_D=0.96$.}
\label{fig:uncertainty-resolution}
\end{figure}

\paragraph{Resolution and operational stability
(Fig.~\ref{fig:uncertainty-resolution}).}
Repeating the sweep at $\Delta x\in\{10,5,2.5\}\,$m shows that the
nominal path-time response remains below $0.43\%$ for every successful
case with $D_{\mathrm{diff}}\le10\,\mathrm{m^2\,s^{-1}}$. At the
finest resolution, minimum clearance increases consistently from
$1.44\,$m at $D_{\mathrm{diff}}=0$ to $1.84\,$m at
$10\,\mathrm{m^2\,s^{-1}}$, whereas the $5$ and $10\,$m grids show
non-monotonic cell-scale variations. The implicit geometric offset is
therefore weak (about $0.4\,$m in this benchmark) and can be obscured by
coarse discretisation; it is not a substitute for the explicit
confidence envelope. Diffusion also shifts the detected empty-domain
wave peak by $3.2$--$6.2\%$ at $D_{\mathrm{diff}}=10$, with the larger
shift on the finer grid. At $N_D=0.96$ the formal explicit bound is
still satisfied, but the $2.5\,$m tracer fails. We consequently use
$N_D\le0.5$ as the empirical operating limit rather than running up to
the formal bound $N_D\le1$.

\begin{figure}[t]
\centering
\includegraphics[width=\linewidth]{fig12_nanshan_uq}
\caption{Zero-wind uncertainty qualification on the $5\,$m Nanshan
geometry ($8.1\times10^7$ cells). (a)~Three routes from station 50 over
the building-footprint map; (b)~building-collision probability from
$4000$ Monte Carlo trajectories per route and diffusivity, with Wilson
$95\%$ intervals; (c)~minimum building clearance minus the local
three-dimensional $r_{95}$. Domain-exit events, including excursions
below ground near the $7.5\,$m-high endpoints, are recorded separately
and are not counted in panel (b).}
\label{fig:uncertainty-nanshan}
\end{figure}

\paragraph{Transfer to a real urban geometry
(Fig.~\ref{fig:uncertainty-nanshan}).}
We next use the full Nanshan mask without a prescribed wind field. One
solve from station 50 serves destinations 17, 19 and 20, separated by
$1.04$, $1.37$ and $2.48\,$km. Their nominal path times are $108.6$,
$132.1$ and $217.6\,$s and their minimum building clearances are
$2.5$, $7.1$ and $12.8\,$m, respectively. At
$D_{\mathrm{diff}}=0.02\,\mathrm{m^2\,s^{-1}}$ the narrowest route
already has a $9.28\%$ building-collision probability (Wilson $95\%$
CI $8.41$--$10.21\%$), while neither of the other routes records a
building collision in $4000$ samples (upper bound $0.10\%$). At
$D_{\mathrm{diff}}=0.5\,\mathrm{m^2\,s^{-1}}$ the three building-risk
estimates are $26.05\%$, $25.23\%$ and $1.33\%$. Thus path duration
alone is not a risk surrogate: local clearance controls whether the
growing uncertainty cloud intersects the urban geometry.

Domain exits are a distinct failure mode. At
$D_{\mathrm{diff}}=0.5\,\mathrm{m^2\,s^{-1}}$ they occur in $24.8$--
$32.5\%$ of samples and are dominated by accumulated vertical
dispersion near the low-altitude endpoints. This exposes the limit of
uncontrolled isotropic Brownian propagation during take-off and
landing: a practical planner must reduce or reset uncertainty using
localisation feedback and a controlled terminal-flight model. We keep
these events separate from building collisions so that the city-geometry
and terminal-control effects are not conflated.

\begin{table}[t]
\centering
\caption{Machine-checkable verification suite (no additional solver
runs required). All checks pass.}
\label{tab:verification}
\small
\begin{tabular}{p{0.34\linewidth}p{0.21\linewidth}p{0.36\linewidth}}
\toprule
Check & Domain & Result \\
\midrule
Reciprocity: forward vs.\ adjoint ($q\to-q$) & synthetic, $5$ wind configs & $0.00\%$ (bitwise); negative control $26\%$ \\
Broadband linearity $\psi_{bb}=\psi_L+\psi_S$ & Nanshan $50\,$m & coverage $100\%$; median $|\Delta T| = 1\,\Delta t$ \\
Multicolor batch $=$ superposed singles & Nanshan $50\,$m, $4$ sources & $0.00\%$ at probe pairs; $3.9\times$ faster \\
Real-domain forward/backward pairs & Nanshan $5\,$m, $48$ pairs & median $0.12\%$ (one boundary outlier $20.5\%$)\footnotemark \\
Wave--dynamics coupled iteration & Nanshan $50\,$m & parameter change $0.396 \to 0.083$ in $5$ iterations \\
\bottomrule
\end{tabular}
\end{table}
\footnotetext{The outlier pair starts and ends adjacent to the
absorbing sponge, consistent with the stated requirement that
first-arrival paths avoid the sponge layer.}

Table~\ref{tab:verification} adds five structural checks that
exercise the identities the efficiency claims of
\S\ref{sec:efficiency} rely on (reciprocity, linearity,
superposition) at machine precision.

% ---------------------------------------------------------------------
\subsection{Comparison with a discrete ground truth on a real city}
\label{sec:groundtruth}

\begin{figure}[t]
\centering
\includegraphics[width=\linewidth]{fig04_groundtruth}
\caption{Wave arrival times versus the discrete ground truth on the
Nanshan domain. (a)~For $555$ station pairs separated by
$700$--$3000\,$m, calibrated wave time is compared with a
$26$-neighbour Dijkstra path shortened by collision-free line-of-sight
segments whenever the continuous metric time decreases.
(b)~Optimality-gap distributions show the old fixed detection
threshold at $5\,$m (grey, long positive tail), the under-resolved
$50\,$m grid (red), and the corrected resolution-scaled threshold at
$5\,$m (blue, compact distribution).}
\label{fig:groundtruth}
\end{figure}

\paragraph{Reference construction.}
The reference is a $26$-neighbour voxel graph over the identical
building mask and station cells, whose edge weights solve the
\emph{same} continuous metric as the PDE coefficients --- horizontal
speed $c_0$ with the regulatory low-altitude band, and the
climb/descent ellipsoid with the vertical offset --- exactly per
edge. Raw $26$-neighbour Dijkstra carries a systematic angular
discretisation error (median $+1.1\%$, up to $+5.0\%$ at $5\,$m), so
all reference paths are shortened by metric-aware string-pulling
(a shortcut is accepted only if it is line-of-sight \emph{and}
reduces metric time); the result is a tight upper bound on the true
optimum, fully attributable as a feasible flight path.

\paragraph{Headline result (Fig.~\ref{fig:groundtruth}).}
Over $555$ valid pairs at $5\,$m the calibrated wave arrival time is
within $\pm 6.5\%$ of the discrete ground truth for $95\%$ of pairs,
with a median of $-4.1\%$. The sign and size of the median offset are
consistent with the bracketing structure: the reference is an upper
bound on the optimum (greedy string-pulling is not globally optimal),
while finite-wavelength diffraction lets the first arrival ``round
corners'' slightly; the true optimum lies between. The detour-ratio
sanity check on the same matrix gives
$T/(d/c_0)$ median $1.12$ / p95 $1.23$ with no pair below $1$: the
free-space lower bound is never violated.

\paragraph{A production bug found by this comparison.}
The first run of this comparison returned a long positive tail
(p95 $36\%$, up to $+230\%$). The cause was \emph{not} the PDE: the
peak-detection threshold, tuned on the $50\,$m grid, was three orders
of magnitude too high at $5\,$m --- source normalisation makes far-field
amplitudes scale as $\Delta x^3$, so weak diffracted first arrivals
in shadow zones fell below threshold and a later, stronger wave
packet was recorded instead ($71$--$86$ of $118$ stations per source).
Scaling the threshold by $(\Delta x / 50\,\mathrm{m})^3$ removed the
tail entirely (Fig.~\ref{fig:groundtruth}b) and was folded back into
the production pipeline. We report this deliberately: the gap
distribution is a property of resolution and detection, not of the
formulation, and the ground-truth harness is what isolates such
regressions.

\begin{table}[t]
\centering
\caption{Wave solver versus the discrete baseline on the $5\,$m
Nanshan domain ($8.1\times10^7$ cells, ${\sim}2\times10^9$ graph
edges).}
\label{tab:comparison}
\small
\begin{tabular}{p{0.24\linewidth}p{0.33\linewidth}p{0.34\linewidth}}
\toprule
 & Wave (ours) & Dijkstra + string-pulling \\
\midrule
One source $\to$ whole domain & $49\,$s (A100) & $70\,$s + $0.3\,$s pull (1 CPU core; graph build ${\sim}9\,$min amortised) \\
Output & full $T$ field, any destination $O(1)$ & per-source distances + predecessors \\
Full $118\times118$ matrix & $118$ solves $=96\,$min; \S\ref{sec:efficiency} reduces per-destination cost to $1$ solve & $118$ searches (parallelisable) \\
Wind / anisotropy / regulation & native PDE coefficients & per-edge closed-form metric (derived here) \\
Memory & ${\sim}1\,$GB (float32 fields) & ${\sim}60\,$GB (CSR graph) \\
Accuracy & $-4.1\%$ median / $6.4\%$ p95 vs.\ reference & $+1.1\%$ angular bias before pulling \\
\bottomrule
\end{tabular}
\end{table}

\begin{table}[t]
\centering
\caption{CUDA fast-marching baseline on Nanshan. Times are JIT-warm
wall-clock measurements on one RTX~5090 and are not compared directly
with the A100 wave timings in Table~\ref{tab:comparison}.}
\label{tab:fmm-baseline}
\small
\begin{tabular}{lll}
\toprule
 & $50\,$m full matrix & $5\,$m single source (station 50) \\
\midrule
Grid cells & $2.87\times10^5$ & $8.08\times10^7$ \\
Scope & 118 sources, 5522 valid pairs & 1 source, 118 endpoints \\
Finite endpoints & $100\%$ & $100\%$ \\
Steady single-source time & $7.83\,$ms & $0.755\,$s \\
CuPy memory-pool peak & $0.84\,$GB & $25.97\,$GB \\
Reference agreement & $+1.37\%$ median; $6.57\%$ p95 absolute & $1.30\%$ median; $2.98\%$ max\footnotemark \\
\bottomrule
\end{tabular}
\end{table}
\footnotetext{The $50\,$m values are against Dijkstra plus
metric-aware string-pulling over the common 700--3000 m validity band.
The $5\,$m values are three pre-registered endpoints (stations 17, 19
and 20) against the calibrated zero-diffusion wave arrivals from
Fig.~\ref{fig:uncertainty-nanshan}; they are an anchor check, not a full-matrix
accuracy estimate.}

\paragraph{Fast marching (Table~\ref{tab:fmm-baseline}).}
For a continuous-field baseline we convert the same horizontal
regulation, $5\,$m/s vertical control speed and $+1\,$m/s vertical
drift exactly to a three-dimensional Randers metric and solve it with
AGD's CUDA \texttt{Rander3} implementation. The algebraic conversion
matches the closed-form edge travel time to $4.2\times10^{-16}$ in a
1000-vector numerical check. On the complete $50\,$m station matrix,
all endpoints are finite and FMM differs from the feasible
string-pulled reference by $+1.37\%$ at the median and $6.57\%$ at the
p95 absolute level. The two directions of each pair agree to within
$0.16\%$ at p95, as expected here because the vertical Randers
one-form integrates to zero between equal-altitude endpoints.

The $5\,$m solve demonstrates that the 5090 can hold the full
$1333\times1443\times42$ problem, but with little margin: the CuPy
pool reaches $25.97\,$GB. After one allocation/warm-up pass, four
repeats take $0.743$--$0.770\,$s (median $0.755\,$s). This is a strong
result for FMM, not a speed claim for the wave solver: in a frozen,
deterministic single-source metric, CUDA fast marching is the more
efficient tool. Table~\ref{tab:comparison} still shows that an
explicit graph is costly, while the distinctive benefits tested for
the wave formulation below are linear source superposition,
transient-field physics and its extension to diffusion-based
uncertainty---not superiority over specialised FMM on its native
static problem.

% ---------------------------------------------------------------------
\subsection{Efficiency: adjoint, superposition, throughput}
\label{sec:efficiency}

\begin{figure}[t]
\centering
\includegraphics[width=\linewidth]{fig05_pipelines}
\caption{Efficiency measurements on the $5\,$m domain.
(a)~Crosstalk of batched (superposed-source) solves versus source
separation: blue gives the maximum pair-time error, orange the
uncontaminated-pair fraction, and the dashed line marks the $6\,$km
safe threshold at twice the pairing radius. (b)~Full-matrix wall time
versus graph-colouring separation; numbers above bars are solve counts
and hatching marks configurations rejected by panel~(a). (c)~Measured
end-to-end cost: multicolour batching reduces the 118-source matrix
from 118 to 82 solves at $0.4\%$ maximum error, while one adjoint solve
replaces 118 forward solves for a destination column at median
$0.35\%$ and p95 $1.0\%$ disagreement.}
\label{fig:pipelines}
\end{figure}

\paragraph{Adjoint (one solve per destination).}
By the reciprocity of Section~2 (flipping the full advection field
$q \to -q$, including the vertical climb/descent offset), one
backward solve from a destination yields the travel time \emph{from
every origin}. Validated against the forward production matrix on
$10$ destinations ($476$ valid pairs), the agreement is median
$0.35\%$, p95 $1.0\%$\footnote{One pair of $476$ disagrees by
$146\%$; on this station pair the production \emph{forward} matrix
itself fails to detect a first peak in the reverse direction --- a
detector-fragility of that geometry, not a reciprocity failure.}.
The wall-clock consequence on the $118$-station district: a
destination column costs $96\,$min as $118$ forward solves and
$49\,$s as one adjoint solve --- a $118\times$ reduction at
sub-percent error (Fig.~\ref{fig:pipelines}c).

\paragraph{Multicolor superposition --- with its measured limits.}
Linearity allows several sources in one solve; the first-arrival
field then serves each source within its own service disk. We
measured the crosstalk of batched solves against production
single-source rows as a function of source separation
(Fig.~\ref{fig:pipelines}a): contamination is severe below
$6\,$km and drops to $0.4\%$ at $6\,$km --- exactly twice the
$3\,$km pairing radius, i.e.\ the separation at which two service
disks can no longer overlap. The criterion is therefore predictable
a priori. Its consequence on \emph{this} domain
($6.7\times7.2\,$km) is modest: a clean colouring of all $118$
stations at the full $3\,$km pairing radius needs $82$ solves
($1.4\times$ saving); for a last-mile matrix with a $1.5\,$km
service radius the same criterion permits $28$ solves
($4.2\times$). The saving scales with the squared ratio of domain
size to service diameter, so batching pays off on large domains and
short-range matrices, and we state the claim in that conditional
form.

\begin{table}[t]
\centering
\caption{Solver throughput on one A100 (production pulse; cold = first
call including compilation with a cleared JIT cache; warm = median of
3).}
\label{tab:throughput}
\begin{tabular}{lrrrrrr}
\toprule
Grid & Cells & $\Delta t$ [s] & Cold [s] & Warm [s] & Steps/s & Cell$\cdot$steps/s \\
\midrule
$50\,$m & $2.9\times10^5$ & $0.1$ & $4.3$ & $0.24$ & $18\,500$ & $5.3\times10^{9}$ \\
$5\,$m & $8.1\times10^7$ & $0.1$ & $55.2$ & $48.9$ & $92$ & $7.4\times10^{12}$ \\
$2.5\,$m subdomain & $1.2\times10^8$ & $0.05$ & $78.0$ & $76.7$ & $59$ & $7.1\times10^{12}$ \\
\bottomrule
\end{tabular}
\end{table}

\paragraph{Throughput (Table~\ref{tab:throughput}).}
Large grids saturate the device at ${\sim}7\times10^{12}$
cell-steps per second (the $5\,$m and $2.5\,$m rows agree), while
the $50\,$m grid is kernel-launch bound, three orders of magnitude
below saturation --- which is why the $50\,$m ablation sweeps of
\S\ref{sec:robustness} cost seconds per solve. Compilation overhead
is $1$--$6\,$s and amortises over array jobs through the persistent
JIT cache. The full $118\times118$ matrix at $5\,$m costs
$96$ A100-minutes ($49\,$s per solve); an eight-GPU array delivers
it in ${\sim}12\,$min wall-clock.

% ---------------------------------------------------------------------
\subsection{Robustness and operating envelope}
\label{sec:robustness}

\begin{figure}[t]
\centering
\includegraphics[width=0.8\linewidth]{fig06_wavelength}
\caption{Wavelength operating point: optimality gap versus dominant
wavelength $\lambda \approx 2\pi c_0 \tau$ (bottom axis) and
$\lambda$ in units of the median street width (top axis), with the
production point and the two-band (broadband) proxy marked.}
\label{fig:wavelength}
\end{figure}

\begin{figure}[t]
\centering
\includegraphics[width=\linewidth]{fig07_wind}
\caption{Spatially varying wind on the $50\,$m domain.
(a)~Path drift under a solenoidal vortex field at $M=0.3$;
(b)~per-pair wind-induced time change $\Delta T$, wave versus
wind-aware Dijkstra; (c)~adjoint reciprocity error versus RMS
$\nabla\!\cdot q$ at fixed $|q|_{\max}$.}
\label{fig:wind}
\end{figure}

\paragraph{Wavelength operating point (Fig.~\ref{fig:wavelength}).}
Sweeping the pulse width over $\tau \in [0.1, 3.2]\,$s
($\lambda \approx 9$--$300\,$m) against the ground truth of
\S\ref{sec:groundtruth} shows three regimes. For
$\lambda \le 0.6\,W$ (where $W = 120\,$m is the median free street
width at cruise altitude) the gap is wavelength-independent ---
median ${\approx}{-4}\%$, p95 $6.5\%$ --- so the error there is a
resolution/detection property, not a wavelength property, and the
production point sits in the middle of this plateau with margin on
both sides. Beyond $\lambda \gtrsim 1.2\,W$ the negative offset
grows (median $-6.4\%$, p95 $9.7\%$ at $\lambda = 2.5\,W$):
diffraction rounds corners more than a point vehicle can follow ---
the measured long-wavelength limit. At the short end the grid
low-pass saturates the effective wavelength near $5\Delta x$, so
wavelengths below ${\sim}30\,$m are unrealisable at $5\,$m
resolution. A two-band combination
$\min(T_{0.2}, T_{1.6})$ (same cost as two solves) restores
long-band timing to plateau accuracy, which is the quantitative
basis for the broadband guidance of Section~2.

\paragraph{Spatially varying wind (Fig.~\ref{fig:wind}).}
On the $50\,$m domain we impose synthetic fields seen identically by
both methods: a solenoidal vortex array (stream function,
$3\,$km cells) and an irrotational compressive checkerboard
(potential flow, $1\,$km period), blended at fixed
$|q|_{\max} = 4.5\,$m/s so that divergence is varied without
confounding wind speed. The reference generalises the Dijkstra
metric to the full wind vector by solving the air-speed ellipsoid
quadratic per edge. Two results. (i)~\emph{Forward physics:} the
per-pair wind-induced change $\Delta T = T(q) - T(0)$ --- a
difference that cancels static discretisation error --- correlates
at $0.86$ with the wind-aware ground truth at $M = 0.3$, with
residuals within the $50\,$m resolution noise; optimal routes drift
downwind as expected (Fig.~\ref{fig:wind}a). The regulatory
low-altitude band under strong wind becomes locally ``supersonic''
(wind exceeds the speed cap); the PDE handles this natively
(perturbations propagate only downwind) and the graph drops the
corresponding directed edges --- the two formulations agree that
such corridors are one-way. (ii)~\emph{Adjoint degradation:}
reciprocity holds exactly only for divergence-free $q$. Measured
forward/backward disagreement rises from a $1.5\%$ discrete floor
(uniform wind) to median $3.7\%$ / p95 $8.5\%$ for the purely
compressive field at RMS $\nabla\!\cdot q = 2.8\times10^{-2}\,$s$^{-1}$
--- a graceful degradation, and urban boundary-layer flows are
vorticity-dominated with divergence an order of magnitude below this
stress test, so the adjoint shortcut remains in its floor regime in
practice.

Two closure experiments --- a broadband tracing ablation on the
hardest long-wavelength pairs, and the wind-inversion and
wave--dynamics coupling loops --- are reported in the appendix; their
outcome in one sentence each: path-tracing robustness, not arrival
timing, is the binding constraint on hard pairs (smoothing the $T$
field before gradient descent lifts success from $10$--$14/20$ to
$17/20$ for every band, with no additional gain from broadband); and
the estimation$\to$planning loop closes (re-planning with wind
recovered from $1254$ noisy track segments halves the
route-time error versus ignoring wind, median $8.5\% \to 3.8\%$),
while the wave--dynamics fixed-point iteration converges in a median
of $4$ iterations from varied initialisations.

% ---------------------------------------------------------------------
\subsection{Concurrent UAVs and dynamic velocity}
\label{sec:dynamic-operations}

\begin{figure}[t]
\centering
\includegraphics[width=\linewidth]{fig13_dynamic_operations}
\caption{Operational stress tests. (a)~Twenty closed-loop missions on
the $50\,$m Nanshan geometry, formed from ten reciprocal station pairs;
each colour denotes the two directions of one obstacle-aware route.
(b)~Number of UAV pairs whose realised separation falls below $50\,$m
when every mission launches at $t=0$ and after delay-only CBS
coordination; the black curve gives the coordinated mean launch delay.
(c)~P95 absolute arrival-time gap against the analytic solution in an
empty $600\,$m cube with a spatially uniform velocity of magnitude
$4.5\,$m/s rotating with the indicated period. Results are shown for
the time-varying solver, a velocity frozen at emission, and an ablation
that omits $\partial_t q$.}
\label{fig:dynamic-operations}
\end{figure}

Figure~\ref{fig:dynamic-operations}(a,b) tests whether nominal routes
remain usable when flown concurrently. We deliberately choose
reciprocal station pairs, an adversarial zero-wind case in which every
pair attempts to traverse the same route head-on. Simultaneous launch
produces $5$, $10$, and $28$ conflicting UAV pairs for fleets of $10$,
$20$, and $40$ missions. Delay-only CBS accepts every mission and
reduces this count to zero in all three cases. The smallest realised
separation over the coordinated runs is $119.4\,$m, above the
$50\,$m threshold; the mean launch delays are $68.5$, $68.5$, and
$70.2\,$s, respectively. These are closed-loop point-mass simulations,
not reservation-only checks: the P95 tracking error remains below
$1.60\,$m and no simulated sample enters a building. The result
therefore establishes feasibility for this deliberately conflicting
traffic pattern, while the delay and doubled makespan quantify the
price of the conservative delay-only policy.

Figure~\ref{fig:dynamic-operations}(c) isolates temporal variation in
the advection velocity without introducing an unverified atmospheric
dataset. For a spatially uniform rotating field, the exact arrival time
is obtained from
$\lVert x-x_s-\int_{t_0}^{t}q(\tau)\,\mathrm d\tau\rVert=c_0(t-t_0)$.
Across $16$ directions, the calibrated dynamic solver has P95 absolute
gaps of $1.36\%$, $1.46\%$, and $1.68\%$ for rotation periods of
$20$, $40$, and $80\,$s. Freezing $q$ at emission instead gives
$47.03\%$, $30.21\%$, and $15.90\%$. Omitting $\partial_tq$ gives
$1.78\%$, $1.79\%$, and $1.68\%$ in this uniform-flow case. Thus the
implemented time-dependent update tracks the analytic characteristic
to within $1.7\%$ P95, whereas a static-velocity approximation fails
when the field changes appreciably during flight. This is a controlled
dynamic-velocity qualification, not evidence for performance in a
measured turbulent wind field.

% ---------------------------------------------------------------------
\subsection{Scale transfer}
\label{sec:transfer}

\begin{figure}[t]
\centering
\includegraphics[width=\linewidth]{fig09_transfer}
\caption{One kernel, two vehicles. (a)~Delivery multirotor on the
$5\,$m Nanshan grid: first-arrival field at cruise height and the
traced $269\,$s path between stations. (b)~Regional aircraft on the
$80\times80\,$km domain: four trajectories around three anisotropic
weather cells at FL260. The configurations differ only in
$c_0$ ($15 \to 200\,$m/s), $c_z$ ($(6,4) \to (25,20)\,$m/s),
resolution ($5 \to 400\,$m) and domain size.}
\label{fig:transfer}
\end{figure}

Figure~\ref{fig:transfer} runs the identical solver and path-extraction
pipeline at two ends of the urban-airspace spectrum: a $15\,$m/s
delivery multirotor threading $5\,$m street canyons, and a
$200\,$m/s regional aircraft climbing to FL260 and routing around
three weather cells in an $80\,$km domain. No code differs between
the two cases --- only the configuration table shown in the caption.
Obstacles enter as a building mask in one case and as smooth
wave-speed suppressions (Section~2.1) in the other; the same
$-\nabla T$ descent produces street-level paths and en-route
climb profiles respectively. The method is a configuration, not a
codebase per vehicle class.

% =====================================================================
% APPENDIX MATERIAL — move the following into the paper's appendix.
% =====================================================================

\section*{Appendix: robustness closures}
\label{sec:appendix-closures}

\paragraph{Broadband ablation on hard pairs.}
We selected the $20$ valid pairs with the most negative
long-wavelength gap in Fig.~\ref{fig:wavelength} (worst $-8.6\%$) and
traced $-\nabla T$ paths on the short-band ($\tau=0.2\,$s),
long-band ($\tau=1.6\,$s) and per-cell-earliest broadband fields,
scoring success as reaching the source within two cells with no
collision along the densely sampled polyline. On the raw $5\,$m
fields, success is $14/20$, $10/20$ and $10/20$ respectively --- but
the failures are tracer stalls, not timing errors (the arrival times
at these pairs are accurate per Fig.~\ref{fig:wavelength}). Smoothing
the finite region of the field by normalised convolution
($\sigma = 2$ cells; buildings keep a sharp repulsive plateau ---
na\"ively smoothing the plateau bleeds hundreds of seconds into
$40\,$m streets and destroys the gradient) lifts all three bands to
$17/20$, with the three remaining failures shared by every band
(a geometric tracing pathology of one source). Broadband therefore
adds no \emph{path-level} success at this resolution; its value is at
the field level --- one two-band excitation yields smooth guidance
and accurate timing simultaneously, and repairs the long-band gap of
Fig.~\ref{fig:wavelength} --- and we state the claim in that narrowed
form.

\paragraph{Wind-inversion closure.}
With the $M=0.3$ vortex field of Fig.~\ref{fig:wind} as truth, we
synthesised flight tracks (twelve sources, five destinations each),
segmented them at $50\,$m, added $0.3\,$m/s observation noise
($1254$ segments), and recovered the wind by Gaussian-kernel
regression ($L = 350\,$m; the kernel attenuates a wavelength
$\Lambda$ by $e^{-k^2L^2/2}$, so $L$ must be small against the
$3\,$km flow structure). The pointwise field recovery is mediocre ---
relative RMS ${\approx}0.73$ --- yet re-solving all sources with the
recovered field halves the route-time error against truth: median
$|\Delta T|/T = 3.8\%$ / p95 $12.9\%$, versus $8.5\%$ / $20.9\%$ when
wind is ignored. Tracks sample exactly the corridors that routing
revisits, so accuracy concentrates where it matters; the
estimation$\to$planning loop of Section~2.6 closes under realistic
sparsity.

\paragraph{Coupled wave--dynamics iteration.}
The fixed-point iteration between the wave parameters
$(c_0, c_z^{\uparrow}, c_z^{\downarrow})$ and a surrogate vehicle
dynamics was run on ten station pairs from five initialisations each
(relaxation $\alpha = 0.5$, $2\%$ tolerance). $33/50$ runs converge,
median four iterations; for seven of ten pairs at least three of five
initialisations converge to the same point within $1$--$2\,$m/s ---
a well-defined attractor --- while two pairs oscillate from every
initialisation. The fixed point itself ($c_0^\ast \approx 5.7\,$m/s)
is a property of the surrogate dynamics evaluated on unsmoothed voxel
paths, whose stair-stepping biases observed speeds low; the
production parameters are regulatory values, not the fixed point of
any simulated dynamics. The experiment therefore supports the
\emph{contraction} property claimed in Section~2 --- fast,
initialisation-robust convergence --- and we make no claim of
parameter recovery.
